% ============================================================
%  Chapitre 6 — Réalisation et Tests
% ============================================================
\chapter{Réalisation et Tests}
\thispagestyle{fancy}

\section{Environnement de développement}

\begin{table}[H]
\centering
\caption{Stack technologique — Smart Farm AI v3.0}
\label{tab:stack}
\begin{tabular}{lll}
\toprule
\textbf{Composant} & \textbf{Technologie} & \textbf{Version} \\
\midrule
\multicolumn{3}{l}{\textbf{Backend}} \\
Langage                  & Python            & 3.11 \\
Framework API            & FastAPI           & 0.100+ \\
ORM                      & SQLAlchemy        & 2.0 \\
Base de données          & PostgreSQL        & 15 / SQLite 3 \\
Authentification         & Python-Jose (JWT) & 3.3.0 \\
\midrule
\multicolumn{3}{l}{\textbf{Intelligence Artificielle}} \\
Vision par ordinateur    & Ultralytics YOLO  & 8.1 \\
LLM local                & Ollama            & 0.1.x \\
LLM Darija               & Labess-7B         & 7B parameters \\
Vecteurs RAG             & ChromaDB          & 0.4+ \\
Embeddings               & sentence-transformers & 2.x \\
Anomalie IoT             & scikit-learn      & 1.4 \\
\midrule
\multicolumn{3}{l}{\textbf{MLOps}} \\
Versionnement données    & DVC               & 3.x \\
Tracking expériences     & MLflow            & 2.x \\
\midrule
\multicolumn{3}{l}{\textbf{Frontend}} \\
Framework UI             & React             & 18.2 \\
Build tool               & Vite              & 5.1 \\
Cartographie             & Leaflet           & 1.9 \\
Graphiques               & Recharts          & 2.x \\
Base offline             & Dexie.js          & 3.x \\
\midrule
\multicolumn{3}{l}{\textbf{IoT}} \\
Microcontrôleur          & ESP32             & -- \\
Framework firmware       & PlatformIO        & 6.x \\
Simulateur               & Wokwi             & -- \\
Protocole communication  & HTTP/JSON         & REST \\
\midrule
\multicolumn{3}{l}{\textbf{Déploiement}} \\
Conteneurisation         & Docker            & 24.x \\
Orchestration            & Docker Compose    & 2.x \\
Cloud                    & Oracle Cloud OCI  & Always Free ARM64 \\
Reverse proxy            & Caddy             & 2.x \\
\bottomrule
\end{tabular}
\end{table}

\section{Endpoints API — Tableau des 20 endpoints principaux}

\begin{longtable}{llp{4.5cm}cc}
\caption{Tableau des endpoints FastAPI principaux} \label{tab:endpoints} \\
\toprule
\textbf{Méthode} & \textbf{URL} & \textbf{Description} & \textbf{Auth} & \textbf{Sprint} \\
\midrule
\endfirsthead
\toprule
\textbf{Méthode} & \textbf{URL} & \textbf{Description} & \textbf{Auth} & \textbf{Sprint} \\
\midrule
\endhead
\bottomrule
\endfoot
POST & \api{/auth/register}         & Création compte + hash bcrypt     & Public & S1 \\
POST & \api{/auth/login}            & Connexion JWT (access + refresh)  & Public & S1 \\
POST & \api{/auth/otp/send}         & Envoi OTP email (6 chiffres)      & Public & S1 \\
POST & \api{/auth/otp/verify}       & Validation OTP 10 minutes         & Public & S1 \\
GET  & \api{/farms/}                & Liste fermes du owner             & JWT   & S1 \\
POST & \api{/farms/}                & Création nouvelle ferme           & JWT   & S1 \\
POST & \api{/iot/telemetry}         & Réception métrique ESP32          & Public & S2 \\
GET  & \api{/iot/latest}            & Dernières métriques NODE\_A/B     & JWT   & S2 \\
GET  & \api{/iot/history}           & Historique CSV avec filtre        & JWT   & S2 \\
WS   & \api{/iot/ws/alerts}         & WebSocket alertes temps réel      & JWT   & S2 \\
GET  & \api{/bee/yards}             & Liste apiaires                    & JWT   & S3 \\
POST & \api{/bee/hives}             & Créer ruche + QR Code             & JWT   & S3 \\
POST & \api{/bee/visits}            & Enregistrer visite ruche          & JWT   & S3 \\
POST & \api{/bee/harvests}          & Enregistrer récolte miel          & JWT   & S3 \\
GET  & \api{/bee/analytics}         & Statistiques apicoles             & JWT   & S3 \\
POST & \api{/cv/scan}               & Inférence YOLO (upload image)     & JWT   & S4 \\
GET  & \api{/cv/models}             & Liste des 8 modèles disponibles   & JWT   & S4 \\
GET  & \api{/cv/history}            & Historique des scans              & JWT   & S4 \\
POST & \api{/assistant/chat}        & Chat Labess-7B + RAG              & JWT   & S5 \\
POST & \api{/assistant/tts}         & Text-to-Speech                    & JWT   & S5 \\
GET  & \api{/mlops/metrics}         & Métriques MLflow dernière run     & JWT   & S6 \\
GET  & \api{/workers/}              & Liste ouvriers farm               & JWT   & S6 \\
\bottomrule
\end{longtable}

\section{Interfaces utilisateur}

\subsection{Tableau de bord Owner (Dashboard)}

Le tableau de bord principal présente en temps réel :

\begin{itemize}
  \item \textbf{Bande KPI} (4 métriques) : Humidité sol, Température ruche, Poids ruche,
        Statut pompe d'irrigation.
  \item \textbf{Widget météo} : Données OpenMeteo (température, précipitations, vent,
        UV) géolocalisées par ferme.
  \item \textbf{Panneau IoT Live} : Dernières lectures NODE\_A et NODE\_B avec
        indicateurs de statut (vert/orange/rouge).
  \item \textbf{6 Badges Espèces} : Accès rapide aux 6 groupes de modèles YOLO
        (apiculture, bétail, cultures, olivier, agrumes, sécurité).
  \item \textbf{Historique Alerts} : Flux des alertes SWARM\_RISK et THERMAL\_FAULT.
\end{itemize}

\subsection{HiveIntel — 7 onglets apiculture}

Le module apiculture HiveIntel est organisé en 7 onglets fonctionnels :

\begin{table}[H]
\centering
\caption{Sous-modules HiveIntel}
\label{tab:hiveIntel}
\begin{tabular}{clp{8cm}}
\toprule
\textbf{\#} & \textbf{Onglet} & \textbf{Fonctionnalités} \\
\midrule
1 & Aperçu        & KPIs rucher : nb ruches, récolte totale, température moy., dernière visite \\
2 & Ruches        & CRUD ruches, QR Codes, état de santé (Excellente/Normale/Attention/Critique) \\
3 & Visites       & Calendrier des visites, observations, photos, noter le comportement \\
4 & Récoltes      & Enregistrement kilos miel, qualité, date, ruche source \\
5 & Stock         & Gestion hausses, état des cadres, rotation du matériel \\
6 & Dépenses      & Suivi budget : traitements Varroa, nourrissement, matériel \\
7 & Télémétrie    & Graphiques NODE\_B : poids (24h/7j/30j), température ruche, humidité ext. \\
\bottomrule
\end{tabular}
\end{table}

\subsection{Scanner Worker PWA}

L'interface Scanner, accessible depuis l'application Worker PWA, propose :

\begin{itemize}
  \item \textbf{13 modèles de détection} organisés en 3 groupes :
        \begin{itemize}
          \item Groupe Plantes : Leaf Disease, Olive, Lemon, Orange, Insects (5 modèles)
          \item Groupe Élevage : Bee OBB, Livestock, Sheep/Goat (3 modèles)
          \item Groupe Sécurité : Fire/Smoke (1 modèle)
        \end{itemize}
  \item \textbf{Caméra live} : inférence en temps réel (streaming WebRTC).
  \item \textbf{Upload image} : analyse d'une photo depuis la galerie.
  \item \textbf{Expert FAB} : bouton flottant pour consulter l'assistant spécialisé
        (profil automatiquement sélectionné selon le modèle actif).
  \item \textbf{Mode offline} : derniers résultats mis en cache via IndexedDB (Dexie.js).
\end{itemize}

\subsection{Assistant Souverain}

\begin{itemize}
  \item Interface de chat type messenger avec historique de conversation.
  \item \textbf{STT} (Speech-to-Text) via WebSpeech API avec prise en charge de l'arabe.
  \item \textbf{TTS} (Text-to-Speech) avec voix arabe sélectionnable.
  \item \textbf{Upload image} pour analyse multimodale via LLaVA.
  \item Indicateur de mode (Local Ollama / Groq API / Lite Mode).
\end{itemize}

\subsection{Télémétrie IoT — Analyse}

La page d'analyse télémétriques présente :
\begin{itemize}
  \item Graphiques Recharts \code{LineChart} : NODE\_A (sol, pression, débit) et
        NODE\_B (poids, température ruche, humidité) sur fenêtres 1h/24h/7j.
  \item Tableau des dernières métriques avec horodatage.
  \item Indicateur de statut connexion IoT (online/offline).
  \item Marqueurs d'anomalies détectées (points rouges sur les courbes).
\end{itemize}

\subsection{Map Center SIG}

\begin{itemize}
  \item Carte Leaflet centrée sur les fermes de l'owner.
  \item Marqueurs personnalisés par type d'exploitation (cultures, élevage, apiculture).
  \item Découverte des vétérinaires et fournisseurs agricoles via API Overpass/OSM.
  \item Calcul de distance et itinéraire Google Maps intégré.
\end{itemize}

\section{Simulations IoT — Wokwi}

Les deux nœuds ESP32 ont été simulés avec Wokwi, permettant de valider le firmware
et le protocole de communication avant le déploiement sur matériel réel.

\begin{lstlisting}[style=jsonstyle, caption={diagram.json NODE\_B (Wokwi)}, label={lst:wokwi}]
{
  "version": 1,
  "author": "FARM AI",
  "editor": "wokwi",
  "parts": [
    { "type": "board-esp32-devkit-v1", "id": "esp32" },
    { "type": "wokwi-dht22", "id": "dht1",
      "attrs": { "pin": "14" } },
    { "type": "wokwi-ds18b20", "id": "ds1",
      "attrs": { "pin": "15" } },
    { "type": "wokwi-resistor", "id": "r1",
      "attrs": { "value": "4700" } },
    { "type": "wokwi-led", "id": "led_alert",
      "attrs": { "color": "red",   "pin": "2" } },
    { "type": "wokwi-led", "id": "led_honey",
      "attrs": { "color": "green", "pin": "4" } }
  ]
}
\end{lstlisting}

\section{Tests}

\subsection{Tests backend (pytest)}

\begin{table}[H]
\centering
\caption{Résultats des tests backend}
\label{tab:tests}
\begin{tabular}{llcc}
\toprule
\textbf{Suite} & \textbf{Cas de test} & \textbf{Passed} & \textbf{Coverage} \\
\midrule
Auth           & Register, Login, OTP, JWT refresh    & 12/12 & 87\% \\
Fermes         & CRUD fermes, isolation multi-tenant  & 8/8   & 92\% \\
IoT Télémétrie & POST métrique, GET latest, anomalie  & 6/6   & 85\% \\
Apiculture     & CRUD ruches, visites, récoltes       & 10/10 & 89\% \\
CV Scan        & Inférence 3 modèles (mock images)    & 6/6   & 78\% \\
\midrule
\textbf{Total} & & \textbf{42/42} & \textbf{86\%} \\
\bottomrule
\end{tabular}
\end{table}

\subsection{Performance YOLO}

Des tests de performance ont été conduits sur CPU Intel Core i7 :

\begin{itemize}
  \item \textbf{Latence inférence moyenne :} 167 ms (p50), 198 ms (p95)
  \item \textbf{Objectif SLA :} $<$ 200 ms ✔
  \item \textbf{Modèle le plus lent :} leaf\_disease (192 ms — 12 classes)
  \item \textbf{Modèle le plus rapide :} fire\_smoke (124 ms — 2 classes)
\end{itemize}

\subsection{Tests IoT — Simulation Wokwi}

La simulation Wokwi a permis de valider :
\begin{itemize}
  \item Envoi HTTP POST toutes les 10 secondes — \textbf{100\% des métriques reçues}
        côté backend FastAPI.
  \item Déclenchement des alertes \code{THERMAL\_FAULT} lorsque la température
        simulée dépasse 38°C.
  \item Reconnexion WiFi automatique après perte de connexion simulée.
  \item Passage en mode \code{Serial-only} (offline) lorsque le WiFi est indisponible.
\end{itemize}
